% =====================================================================
%  CONJURA WORKING DOCUMENT
%  Parallel presampling: the conjecture, the transfer to the ordinary
%  bit-fixing model, and multi-source extraction from it. ell stages.
%  Supersedes parallel-ai-bf.tex and extraction-from-parallel.tex.
%  Requires conjura-conjecture.cls in the same directory.
% =====================================================================
\documentclass{conjura-conjecture}

\runninghead{PARALLEL PRESAMPLING}

\newcommand{\Fun}{\mathsf{Fun}}
\newcommand{\Adv}{\mathsf{Adv}}
\newcommand{\Succ}{\mathrm{Succ}}
\newcommand{\Prob}[1]{\Pr\!\left[#1\right]}
\newcommand{\Exp}{\mathbb{E}}
\newcommand{\bits}{\{0,1\}}
\newcommand{\vz}{\mathbf{z}}
\newcommand{\vr}{\mathbf{r}}
\newcommand{\vx}{\mathbf{x}}
\newcommand{\vl}{\boldsymbol{\ell}}
\newcommand{\SAI}{\mathsf{SAI\text{-}RO}}
\newcommand{\SBF}{\mathsf{SBF\text{-}RO}}
\newcommand{\BF}{\mathsf{BF\text{-}RO}}
\newcommand{\pre}{\mathsf{pre}}
\newcommand{\main}{\mathsf{main}}
\newcommand{\Hmin}{H_{\infty}}
\newcommand{\Real}{\mathsf{Real}}
\newcommand{\Dec}{\mathsf{Dec}}
\newcommand{\Tcomb}{T^{\mathsf{comb}}_{G}}
\newcommand{\merge}{\mathbin{\triangleright}}
\newcommand{\Ext}{\mathsf{Ext}}

\begin{document}

\cjkicker{CONJURA WORKING DOCUMENT}
\cjtitle{Parallel Presampling}
\cjsubtitle{The conjecture, the transfer to the ordinary bit-fixing
  model, and multi-source extraction}
\cjstatus{The conjecture of Section~\ref{sec:lemma} is open at every
  $\ell \ge 2$ and is stated on its own in \texttt{statement.tex} beside
  this file, which is the canonical statement; this document is the
  development. Everything else here is proved from it. The case of
  deterministic preprocessing is immediate from~\cite{CDGS}
  (Remark~\ref{rem:det}), so the whole content lies in the private coins.
  Theorem~\ref{thm:ext} is not a new statement:~\cite[Thm.~3]{CFHS}
  proves it unconditionally by compression, with an incomparable bound
  (Remark~\ref{rem:priorart}); what is derived here is the route.}
\cjcategory{\emph{Idealized Models \& Non-Uniformity}.}

\vspace{0.4em}

\begin{informalconjecture}
Several parties read an entire random function, without communicating and
on independent coins, and each writes down a short string. An online
adversary is then given all the strings and a bounded number of queries to
the function. The conjecture says that each party can, on its own, mark a
short list of points of the function, so that the online adversary cannot
tell the real function from one fixed in advance on all the lists together
and fresh everywhere else. Because the parties mark points of the
\emph{same} function, no two lists disagree where they overlap, so no
arbitration is needed. Two things follow: any bound already proved in the
ordinary bit-fixing model applies verbatim against parallel preprocessing,
and a random oracle extracts randomness from any number of independent
unpredictable sources.
\end{informalconjecture}

\section{The setting}\label{sec:setting}

Presampling is the standard route to non-uniform security in idealized
models. A random oracle together with $S$ bits of oracle-dependent advice
is close to a mixture of oracles fixed on $P$ points and uniform
elsewhere, at a cost of roughly $ST/P$ against $T$-query
adversaries~\cite{Unruh,DGK,CDGS}. The technique assumes a \emph{single}
preprocessing stage, and the reduction carrying a bit-fixing bound back to
the auxiliary-input model~\cite[Thm.~5]{CDGS} exploits that assumption
directly: the bit-fixing attacker simulates the one preprocessor
internally, on a table it samples itself, and can therefore inspect the
whole of that table together with the whole of the advice before deciding
which points to fix.

This document asks for the analogue with $\ell \ge 2$ preprocessing stages
running in parallel: they read the same oracle, use independent private
coins, and never communicate. The obstruction is not that any one leakage
is long. It is that no bit-fixing preprocessor can see the other stages'
leakage, because that leakage depends on private coins it does not hold,
and the points needing to be fixed are in general a function of all of
them. If the stages were deterministic there would be nothing to prove,
since any one of them could recompute the others' outputs from the oracle
and invoke the single-source reduction (Remark~\ref{rem:det}).
Randomization is what makes the question a question.

Two consequences are developed, and they are what the conjecture is for.
Section~\ref{sec:theorem} carries any bound proved in the \emph{ordinary}
bit-fixing model over to parallel preprocessing, so that existing bounds
transfer without being reproved. Section~\ref{sec:ext} runs that transfer
on one application and obtains multi-source extraction: the value of the
oracle at a tuple of unpredictable points is indistinguishable from
uniform. That is the application of~\cite[\S1.4]{CFHS}, which raises the
decomposition question and reaches its own bounds by compression instead
--- compression being what stalls for non-monolithic constructions such as
Merkle--Damg\aa{}rd and sponge, and hence the reason to want a
decomposition route at all.

Two features of the intended statement are worth naming at the outset,
because both are choices and both could have gone otherwise. First, each
preprocessor is allowed to choose its list of points using the
\emph{entire} oracle, not merely its own leakage; this is what lets it fix
those points to the values the oracle actually takes there, and hence what
makes the lists automatically compatible
(Proposition~\ref{prop:noconflict}). Second, the transfer is stated
against the ordinary bit-fixing model rather than a split one, which is
legitimate because a split bit-fixing attacker collapses into an ordinary
one (Proposition~\ref{prop:merge}) and is what makes it useful. Not every
application can afford that collapse, and extraction is one that cannot
(Remark~\ref{rem:nomerge}).

\section{Notation and parameters}\label{sec:notation}

For $n \in \mathbb{N}$ write $[n] := \{1,\dots,n\}$, and let $\bits^{*}$
denote the set of finite binary strings. All logarithms are base two. For
a finite set $\mathcal{X}$ we write $x \sample \mathcal{X}$ for a uniform
draw, independent of all others unless stated otherwise. For a random
variable $W$, $\Hmin(W) := -\log\max_w \Prob{W = w}$ and $H(W)$ is Shannon
entropy.

An algorithm written $A^{f}$ has adaptive oracle access to $f$ and learns
nothing about $f$ except through queries. Algorithms are computationally
unbounded and may be randomized, so that ``algorithm'' may be read
throughout as ``arbitrary, possibly randomized, function''; only query
counts are ever bounded. An \emph{absolute constant} does not depend on
any parameter below. Indices $i$ and $j$ range over $[\ell]$.

Fix throughout:
\begin{itemize}
\item $\ell \in \mathbb{N}$ with $\ell \ge 2$, the number of parallel
  preprocessing stages.
\item $N, M \in \mathbb{N}$, and $\Fun := \{\, f : [N] \to [M] \,\}$, a
  function being identified with its table.
\item $S_1,\dots,S_\ell \in \mathbb{N}$, the advice lengths, and
  $S := \sum_i S_i$.
\item $T \in \mathbb{N}$, the online query budget.
\item $P_1,\dots,P_\ell \in \mathbb{N}$, the presampling budgets, and
  $P := \sum_i P_i$.
\item $\gamma \in (0,1)$, an additive slack, and
  $A_i := S_i + \log\gamma^{-1}$.
\item $\mathcal{R}_1,\dots,\mathcal{R}_\ell$, the private coin spaces.
\end{itemize}

\section{The models}\label{sec:models}

An oracle $\mathcal{O}$ has two interfaces, $\mathcal{O}.\pre$ and
$\mathcal{O}.\main$, the first accessible only before any call to the
second. We follow~\cite[\S2.2.1]{CDGS} except that $\pre$ is accessed by
$\ell$ parties.

\begin{definition}[the two oracles]\label{def:oracles}
\emph{Split auxiliary-input random oracle} $\SAI(N,M)$: samples
$F \sample \Fun$; outputs $F$ at $\pre$, separately to each of the $\ell$
parties that query it; answers $x \in [N]$ at $\main$ by $F[x]$.

\emph{Split bit-fixing random oracle} $\SBF(P_1,\dots,P_\ell,N,M)$:
samples $F \sample \Fun$; accepts at $\pre$, from party $i$, a list $L_i$
of at most $P_i$ query/answer pairs with distinct queries; answers
$x \in [N]$ at $\main$ by $(L_1 \merge \cdots \merge L_\ell)[x]$ if $x$
lies in the domain of some $L_i$, and by $F[x]$ otherwise, where the
\emph{priority merge} $L_1 \merge \cdots \merge L_\ell$ takes at a point
$u$ the value $L_i[u]$ for the least $i$ whose list covers $u$.

The merge is used only so that $\main$ is defined on every input. Under
Conjecture~\ref{conj:main} no two lists disagree and the choice is
immaterial; Section~\ref{sec:conflict} treats the case where they may.
\end{definition}

\begin{definition}[split attackers]\label{def:attackers}
An \emph{$(S_1,\dots,S_\ell,T,p)$-split attacker} is a tuple
$A = (A^{0}_1,\dots,A^{0}_\ell,A^{1})$ in which
\begin{itemize}
\item $A^{0}_i$ interacts with $\mathcal{O}.\pre$, draws private coins
  $r_i \sample \mathcal{R}_i$, and outputs $z_i \in \bits^{S_i}$; the
  stages run on \emph{mutually independent} coins and do not communicate;
\item $A^{1}$ takes $\vz := (z_1,\dots,z_\ell)$, makes at most $T$ queries
  to $\mathcal{O}.\main$, and meets whatever further restrictions $p$
  imposes.
\end{itemize}
Write $f_i(F;r_i) := z_i$ for the randomized leakage map computed by
$A^{0}_i$. Taking $S_j = 0$ and $A^{0}_j$ trivial for every $j \ge 2$
recovers the $(S,T,p)$-attackers of~\cite[Def.~2]{CDGS}.
\end{definition}

\begin{definition}[applications, advantage, advice restriction]
\label{def:app}
Fix maps $\pi_i : \bits^{S_i} \to \bits^{*}$ and put
$\pi(\vz) := (\pi_1(z_1),\dots,\pi_\ell(z_\ell))$. A
\emph{$\pi$-restricted application} $G$ in the $\mathcal{O}$-model is
given by a challenger $\mathsf{C}$ with access to $\mathcal{O}.\main$ that
receives $\vz$ and interacts with $A^{1}$, which receives $\pi(\vz)$ in
place of $\vz$; $\mathsf{C}$ outputs a bit and
\[
  \Succ_{G,\mathcal{O}}(A) \;:=\;
  \Prob{\bigl(A^{1}\bigr)^{\mathcal{O}.\main}\!\bigl(\pi(\vz)\bigr)
  \leftrightarrow \mathsf{C}^{\mathcal{O}.\main}(\vz) = 1} ,
\]
the probability over every draw involved, including $F$, the coins $\vr$
and those of $A^{1}$ and $\mathsf{C}$. For an
\emph{indistinguishability} application $\Adv_{G,\mathcal{O}}(A) :=
2|\Succ_{G,\mathcal{O}}(A) - \tfrac12|$; for an \emph{unpredictability}
application $\Adv_{G,\mathcal{O}}(A) := \Succ_{G,\mathcal{O}}(A)$. $G$ is
$\bigl((S_1,\dots,S_\ell,T,p),\varepsilon\bigr)$-secure in the
$\mathcal{O}$-model if $\Adv_{G,\mathcal{O}}(A) \le \varepsilon$ for every
$(S_1,\dots,S_\ell,T,p)$-split attacker $A$.

The \emph{combination} of $A$ and $\mathsf{C}$ is the single oracle
algorithm $D$ that takes $\vz$, runs $A^{1}(\pi(\vz)) \leftrightarrow
\mathsf{C}(\vz)$ forwarding all oracle queries, and outputs $\mathsf{C}$'s
bit; $\Tcomb$ denotes an upper bound on its query count, the
\emph{combined query complexity}.

Taking $\pi$ the identity and $\mathsf{C}$ ignoring its input gives the
ordinary notion, in which the advice goes to the online stage and to
nobody else. The restriction is carried from the start because one
application below needs it in both directions at once
(Remark~\ref{rem:whypi}), and because it costs nothing: only the
combination $D$ ever appears in a proof, and $D$ reads $\vz$ either way.
\end{definition}

\section{The parallel decomposition lemma}\label{sec:lemma}

The object the conjecture asks for is a tuple of \emph{selectors}, one per
preprocessing stage, each reading the whole oracle but only its own
leakage and its own coins.

\begin{definition}[split selector tuples]\label{def:selector}
A \emph{split selector tuple} with budgets $(P_1,\dots,P_\ell)$ is a tuple
of maps
\[
\begin{aligned}
  \Lambda_i &\;:\; \Fun \times \bits^{S_i} \times \mathcal{R}_i
  \;\longrightarrow\; 2^{[N]} \times [M]^{*}, \\
  \Lambda_i(F,z_i,r_i) &\;=\; (I_i, a_i),
\end{aligned}
\]
for $i \in [\ell]$, with $|I_i| \le P_i$ and $a_i : I_i \to [M]$. The
tuple is \emph{consistent} if $a_i = F|_{I_i}$ for every $i$ and all
inputs, that is, if each selector fixes its points to the values the
oracle actually takes there.

Given $F$, $\vz$ and $\vr := (r_1,\dots,r_\ell)$, write
$I := \bigcup_i I_i$ and let $Y[\Lambda]$ denote the random function
taking the value $(a_1 \merge \cdots \merge a_\ell)[u]$ at each $u \in I$
and an independent uniform value of $[M]$ at each $u \notin I$. Once
$(F,\vz,\vr)$ is fixed, $Y[\Lambda]$ is a $P$-bit-fixing source, and
$|I| \le P$ always.
\end{definition}

\begin{conjecture}[parallel decomposition]\label{conj:main}
There is an absolute constant $c$ such that the following holds for all
$\ell \ge 2$, all $N$ and $M$, all randomized leakage maps
$f_1,\dots,f_\ell$ with output lengths $S_1,\dots,S_\ell$ and mutually
independent coins, all $P_1,\dots,P_\ell \in \mathbb{N}$ and all
$\gamma \in (0,1)$. There is a \emph{consistent} split selector tuple
$(\Lambda_1,\dots,\Lambda_\ell)$ with budgets $(P_1,\dots,P_\ell)$,
depending only on $f_1,\dots,f_\ell$, $P_1,\dots,P_\ell$ and $\gamma$,
such that, writing
\[
  \Delta(T) \;:=\; c\,T\sum_{i=1}^{\ell}\frac{A_i}{P_i}
  \;=\; c\,T\sum_{i=1}^{\ell}\frac{S_i + \log\gamma^{-1}}{P_i} ,
\]
the following two bounds hold for every $T \in \mathbb{N}$ and every
distinguisher $D$ that takes a tuple of strings and makes at most $T$
oracle queries. Draw $X \sample \Fun$, then $r_1,\dots,r_\ell$
independently, set $z_i := f_i(X;r_i)$, and draw $Y[\Lambda]$ as in
Definition~\ref{def:selector} \emph{without} resampling $X$, $\vz$ or
$\vr$. Then
\[
  \Bigl|\, \Prob{D^{X}(\vz) = 1} - \Prob{D^{Y[\Lambda]}(\vz) = 1}\,\Bigr|
  \;\le\; \Delta(T) + \gamma ,
\]
and, with $\Delta^{+}(T)$ denoting $\Delta(T)$ with each
$\log\gamma^{-1}$ replaced by $2\log\gamma^{-1}$,
\[
  \Prob{D^{X}(\vz) = 1}
  \;\le\; 2^{\Delta^{+}(T)}\cdot\Prob{D^{Y[\Lambda]}(\vz) = 1} + 2\gamma .
\]
\end{conjecture}

\begin{remark}[reduction to the known case]\label{rem:reduces}
Taking $f_j$ constant, so $S_j = 0$, and $P_j = 1$, for every $j \ge 2$,
the statement is Lemma~1 of~\cite{CDGS} up to the constant $c$, with the
single difference that $f_1$ is permitted to be randomized. That
difference is free: Lemma~\ref{lem:randleak} below shows that the entropy
bookkeeping on which the single-source proof rests survives randomized
leakage verbatim. The entire content of Conjecture~\ref{conj:main} is
therefore that the fixed set may be presented as a union
$I_1 \cup \cdots \cup I_\ell$ whose parts are chosen without reference to
one another's leakage.
\end{remark}

\begin{remark}[quantifier order]\label{rem:order}
The selector tuple is chosen after $P_1,\dots,P_\ell$ and $\gamma$ but
\emph{before} $T$ and before $D$. Relaxing this makes the statement easier
in a way neither consequence can use: both fix their query budget only
after the attacker, and Theorem~\ref{thm:ext} moreover needs one tuple to
serve two distinguishers at two different budgets.
\end{remark}

\begin{remark}[why the selectors read the whole oracle]\label{rem:allofF}
$\Lambda_i$ is given $F$ in full, and not merely $z_i$. This is what makes
consistency achievable: a selector that knows the table can fix its points
to the values they actually take. Consistency in turn is what makes the
lists compatible, which is the whole of
Proposition~\ref{prop:noconflict}. It is also what the single-source
construction already does, where the bit-fixing preprocessor picks the
component of the decomposition containing the table it sampled. The
selectors are \emph{not} given the other stages' leakage or coins, and
that restriction, not the input $F$, is where the difficulty lies.
\end{remark}

\begin{remark}[separate decomposition is not enough, and what replaces it]
\label{rem:separate}
The obvious construction --- run the single-source selector on each $z_i$
separately and take the union --- fails already at $\ell = 2$. Identify
$[M]$ with $\mathbb{Z}_M$, let $f_1$ draw $u \sample [N]$ and output
$z_1 := \bigl(u, \sum_{x \ne u} X(x)\bigr)$, let $f_2$ draw
$v \sample [N]$ and output $z_2 := \bigl(v, \sum_{x \ne v} X(x)\bigr)$,
both sums in $\mathbb{Z}_M$, and take $f_j$ constant for $j \ge 3$. Each
leakage constrains $X$ by a single linear equation on $N-1$ coordinates,
so each conditional source is $(1-\delta)$-dense for $\delta = 1/(N-1)$ ---
the binding subset is $[N]\setminus\{u\}$, on which the min-entropy is
$(N-2)\log M$ --- and the single-source selector fixes nothing at all. But
conditioned on both, and on $u \ne v$, the two equations differ by
$X(v)-X(u)$, so the two-query $D$ that reads the oracle at $u$ and at $v$,
both of which its input names, and compares $X(v)-X(u)$ against the
difference of the leaked sums distinguishes with advantage $1-1/M$ from
any source uniform on $\{u,v\}$. This is the attack
of~\cite[\S1.4]{CFHS}.

The example does not contradict the conjecture; it locates what the
selectors have to do. Setting $I_1 := \{u\}$ and $I_2 := \{v\}$, each
computable from its own leakage alone, leaves both equations saying the
same thing about the remaining coordinates, namely a single linear
constraint on $N-2$ of them, which is $(1-\delta)$-dense for
$\delta = 1/(N-2)$. Each stage fixes the coordinate its \emph{own} coins
point at, and the joint structure takes care of itself. Whether that is a
general phenomenon is exactly what is open.
\end{remark}

\begin{remark}[the shape of the error term is a guess]\label{rem:shape}
The single-source lemma of~\cite{CDGS} fixes $\Delta(T)$ only for
$\ell = 1$. Everything beyond that --- in particular that the $\ell$
stages contribute a \emph{sum} $\sum_i A_i/P_i$, one term per stage, each
of the single-source shape --- is extrapolation. It is the form one gets
by imagining that each stage pays for its own leakage against its own
budget and that nothing is paid for the interaction between them, and it
is what makes Theorem~\ref{thm:aibf} degrade by a factor $\ell$ rather
than worse (Remark~\ref{rem:payoff}). No evidence offered here fixes it. A
true statement with $\Delta(T) = c\,T\,\phi(\ell)\sum_i A_i/P_i$ for some
growing $\phi$, or one in which the budgets enter through $\min_i P_i$
rather than one by one, would leave the structure of every proof below
intact and change every bound they produce. Readers should treat the
constant $c$ as the only thing this document is casual about, and the
functional form as a second thing.
\end{remark}

\section{The parallel reduction}\label{sec:theorem}

Two elementary facts do the work. The first is why no arbitration between
the lists is needed; the second is why the hypothesis may be stated in the
ordinary bit-fixing model, so that no bound has to be reproved.

\begin{proposition}[consistent lists never conflict]\label{prop:noconflict}
Let $(\Lambda_1,\dots,\Lambda_\ell)$ be consistent, and for a table $F$
put $L_i := \{(u,a_i(u)) : u \in I_i\}$. Then $L_i$ and $L_j$ agree at
every point of $I_i \cap I_j$, for every $i$ and $j$, so $\bigcup_i L_i$
is a well-defined list of at most $P$ query/answer pairs and
$L_1 \merge \cdots \merge L_\ell = \bigcup_i L_i$.
\end{proposition}

\begin{proof}
For $u \in I_i \cap I_j$ consistency gives $a_i(u) = F(u) = a_j(u)$. The
size bound is $\bigl|\bigcup_i I_i\bigr| \le \sum_i |I_i| \le \sum_i P_i =
P$.
\end{proof}

The next proposition is the transfer proper, and it is stated at the level
of the \emph{split} bit-fixing model, because that is the level at which
it is true of every application. Only after it does
Proposition~\ref{prop:merge} collapse $\SBF$ into $\BF$, and
Section~\ref{sec:ext} is an application for which that last step is
unavailable.

\begin{proposition}[split auxiliary input to split bit-fixing]
\label{prop:sbf}
Assume Conjecture~\ref{conj:main} with constant $c$. Let $G$ be a
$\pi$-restricted application with combined query complexity $\Tcomb$ and
let $A = (A^{0}_1,\dots,A^{0}_\ell,A^{1})$ be an
$(S_1,\dots,S_\ell,T,p)$-split attacker against it, with leakage maps
$f_1,\dots,f_\ell$, and let $(\Lambda_1,\dots,\Lambda_\ell)$ be the
consistent selector tuple the conjecture supplies for
$f_1,\dots,f_\ell,P_1,\dots,P_\ell,\gamma$. Define
$B = (B^{0}_1,\dots,B^{0}_\ell,B^{1})$ against
$\SBF(P_1,\dots,P_\ell)$: the preprocessing stages share a table
$F \sample \Fun$, stage $B^{0}_i$ draws $r_i \sample \mathcal{R}_i$,
computes $z_i := f_i(F;r_i)$ and $(I_i,a_i) := \Lambda_i(F,z_i,r_i)$,
presets $L_i := \{(u,F(u)) : u \in I_i\}$ and outputs $z_i$; and
$B^{1} := A^{1}$. Then
\[
  \bigl|\Succ_{G,\SAI}(A) - \Succ_{G,\SBF}(B)\bigr|
  \;\le\; \Delta(\Tcomb) + \gamma ,
\]
\[
  \Succ_{G,\SAI}(A)
  \;\le\; 2^{\Delta^{+}(\Tcomb)}\,\Succ_{G,\SBF}(B) + 2\gamma .
\]
\end{proposition}

\begin{proof}
The selector tuple may be used because it depends only on
$f_1,\dots,f_\ell,P_1,\dots,P_\ell,\gamma$, all of which are determined by
$A$, and not on $T$, on $\mathsf{C}$ or on $G$ (Remark~\ref{rem:order}).
$B$ respects the presampling budgets, since $|I_i| \le P_i$, and its online
stage is $A^{1}$ unchanged, so it meets $T$, the restrictions $p$, and the
advice restriction $\pi$, all of which constrain the online stage only.
Its stages share the table $F$, which Definition~\ref{def:attackers} does
not permit; Proposition~\ref{prop:derand} removes the sharing, and
Remark~\ref{rem:derandcost} says when that repair may be used and when it
may not.

\emph{The oracle $B^{1}$ sees is $Y[\Lambda]$.} By
Proposition~\ref{prop:noconflict} the lists agree where they meet, so the
preset is $\bigcup_i L_i$, of size at most $P$, and it fixes the oracle to
$F$'s own values on $I$. Off $I$ the $\SBF$ oracle answers with its own
internally sampled table, uniform, independent across points and
independent of $F$, since $F$ enters the experiment only through $\vz$ and
the lists. Hence, jointly with $\vz$, the oracle offered at $\main$ is
distributed exactly as $Y[\Lambda]$ with $X := F$. In the $\SAI$-model, by
contrast, the oracle at $\main$ is the same table the preprocessors read,
so it is distributed exactly as $X$ with $z_i = f_i(X;r_i)$.

\emph{Applying the conjecture.} Let $D$ be the combination of $A^{1}$ and
$\mathsf{C}$: it takes $\vz$, makes at most $\Tcomb$ oracle queries and
outputs a bit. That the challenger reads $\vz$ and the online stage reads
only $\pi(\vz)$ is immaterial here, and is the only point at which the
advice restriction could have entered: $D$ is a single oracle algorithm
taking the tuple $\vz$, which is all Conjecture~\ref{conj:main} requires of
a distinguisher, and whether $\vz$ reaches it through one part or the other
does not change the object. The two paragraphs above identify
\[
\begin{aligned}
  \Succ_{G,\SAI}(A) &= \Prob{D^{X}(\vz) = 1}, \\
  \Succ_{G,\SBF}(B) &= \Prob{D^{Y[\Lambda]}(\vz) = 1} ,
\end{aligned}
\]
and the two bounds of the conjecture at $T := \Tcomb$ are the two claims.
\end{proof}

\begin{proposition}[a split bit-fixing attacker is a bit-fixing attacker]
\label{prop:merge}
Let $B$ be an $(S_1,\dots,S_\ell,T,p)$-split attacker in the
$\SBF(P_1,\dots,P_\ell)$-model whose preprocessing stages may share
randomness. There is an $(S,T,p)$-attacker $\widehat{B}$ in the
$\BF(P)$-model of~\cite[\S2.2.1]{CDGS} with
$\Succ_{G,\BF}(\widehat{B}) = \Succ_{G,\SBF}(B)$ for every application $G$.
\end{proposition}

\begin{proof}
$\widehat{B}^{0}$ draws the shared randomness and all private coin
strings, runs $B^{0}_1,\dots,B^{0}_\ell$ internally, presets the oracle
with $L_1 \merge \cdots \merge L_\ell$, which has at most $P$ pairs, and
outputs $z := z_1\|\cdots\|z_\ell$, of length $S$. $\widehat{B}^{1}$
parses $z$ and runs $B^{1}$. All stages are computationally unbounded, so
the internal simulation is free, and every distribution in the two
experiments coincides. The online stage is unchanged, so $p$ and $T$ are
met.
\end{proof}

\begin{theorem}[parallel auxiliary-input to bit-fixing]\label{thm:aibf}
Assume Conjecture~\ref{conj:main}, with constant $c$. Let
$P_1,\dots,P_\ell \in \mathbb{N}$ and $\gamma \in (0,1)$, and let $G$ be a
$\pi$-restricted application with combined query complexity $\Tcomb$.
\begin{enumerate}
\item If $G$ is $\bigl((S,T,p),\varepsilon'\bigr)$-secure in the
  $\BF(P)$-model, then $G$ is
  $\bigl((S_1,\dots,S_\ell,T,p),\varepsilon\bigr)$-secure in the
  $\SAI$-model for
  \[
    \varepsilon \;\le\; \varepsilon'
    \;+\; 2c\,\Tcomb\sum_{i=1}^{\ell}\frac{S_i+\log\gamma^{-1}}{P_i}
    \;+\; 2\gamma .
  \]
\item If moreover $G$ is an unpredictability application and
  $P_i \ge \ell\,c\,(S_i + 2\log\gamma^{-1})\,\Tcomb$ for every $i$, then
  $\varepsilon \le 2\varepsilon' + 2\gamma$.
\end{enumerate}
\end{theorem}

\begin{proof}
Fix a split attacker $A$ and let $B$ be as in
Proposition~\ref{prop:sbf}. For an unpredictability application the
advantage is the success probability, so its first bound gives
$\Adv_{G,\SAI}(A) \le \Adv_{G,\SBF}(B) + \Delta(\Tcomb) + \gamma$; for an
indistinguishability application, $2|s-\tfrac12| \le 2|s'-\tfrac12| +
2|s-s'|$ applied to $s := \Succ_{G,\SAI}(A)$ and $s' := \Succ_{G,\SBF}(B)$
gives $\Adv_{G,\SAI}(A) \le \Adv_{G,\SBF}(B) + 2\Delta(\Tcomb) + 2\gamma$.
The larger of the two covers both. By Proposition~\ref{prop:merge} there
is an $(S,T,p)$-attacker $\widehat{B}$ against $\BF(P)$ with the same
success probability, hence the same advantage, and security in that model
gives $\Adv_{G,\SBF}(B) = \Adv_{G,\BF}(\widehat{B}) \le \varepsilon'$.
That is part~(1).

For part~(2) use the second bound of Proposition~\ref{prop:sbf}. The
hypothesis makes each of the $\ell$ summands of $\Delta^{+}(\Tcomb)$ at
most $1/\ell$, so $\Delta^{+}(\Tcomb) \le 1$ and
$2^{\Delta^{+}(\Tcomb)} \le 2$, whence $\Succ_{G,\SAI}(A) \le
2\Succ_{G,\SBF}(B) + 2\gamma$, which for an unpredictability application
is $\Adv_{G,\SAI}(A) \le 2\varepsilon' + 2\gamma$.
\end{proof}

\begin{proposition}[the shared table is removable]\label{prop:derand}
Let $B$ be the split attacker built in Proposition~\ref{prop:sbf}, whose
stages share the table $F$. There is a split attacker $B^{*}$ whose stages
share nothing, are individually randomized, and satisfy
$\Succ_{G,\SBF}(B^{*}) \ge \Succ_{G,\SBF}(B)$.
\end{proposition}

\begin{proof}
For $F \in \Fun$ let $B_F$ be $B$ with the shared table hardwired into all
stages, so that $B_F$'s stages share nothing and draw only
$r_1,\dots,r_\ell$, independently. Then $\Succ_{G,\SBF}(B) =
\Exp_{F \sample \Fun}[\Succ_{G,\SBF}(B_F)]$, so some $F^{*}$ attains at
least the mean; take $B^{*} := B_{F^{*}}$. Preprocessing stages are
computationally unbounded and non-uniform, so hardwiring a table is free.
\end{proof}

\begin{remark}[hardwiring conditions on a chosen table]\label{rem:derandcost}
Proposition~\ref{prop:derand} is innocuous for a restriction $p$ that
constrains only the online stage, which is the case throughout
Section~\ref{sec:theorem} and in the applications
of~\cite[\S\S3--4]{CDGS}. It is not innocuous in general: the hardwired
$F^{*}$ is chosen to \emph{maximise} success, and the stages' outputs are
thereafter conditioned on that table, so a restriction asserting something
about the \emph{distribution} of the advice --- unpredictability, say ---
need not survive. Section~\ref{sec:ext} is exactly such a case, and uses
the shared-table $B$ directly, which Proposition~\ref{prop:merge} already
tolerates. Nothing there needs $B$ to be admissible: only the
identification of the $\SBF$ oracle with $Y[\Lambda]$ is used, and the
bit-fixing bound is proved on the probability space directly, without
quantifying over attackers.
\end{remark}

\begin{remark}[what the theorem buys]\label{rem:payoff}
The hypothesis of Theorem~\ref{thm:aibf} is security in the
\emph{ordinary} bit-fixing model, so every bound already proved there
transfers to $\ell$ parallel preprocessing stages with no new analysis:
one instantiates $\varepsilon'$ from the existing proof and pays the
additive term. Applied to the applications of~\cite[\S\S3--4]{CDGS} ---
one-way functions, pseudorandom generators and functions,
message-authentication codes, and collision resistance of
Merkle--Damg\aa{}rd --- it yields the corresponding statements against
split advice at $S$ and $P$, the additive term degrading by at most a
factor of $\ell$ relative to the single-stage bound when all
$P_i = P/\ell$ and all $S_i$ are equal. Not every application enjoys this:
one whose security depends on a restriction that the collapse of
Proposition~\ref{prop:merge} destroys must stop at
Proposition~\ref{prop:sbf} and prove its bit-fixing bound in the split
model. Extraction is such an application (Remark~\ref{rem:nomerge}).
\end{remark}

\section{Conflicting lists}\label{sec:conflict}

Consistency is what removes arbitration, and it is worth recording what
happens without it, since a weaker decomposition lemma may be easier to
prove. Call a split selector tuple \emph{free} if the $a_i$ are
unrestricted, and let $\mathsf{Conf}$ be the event that $a_i$ and $a_j$
disagree at some point of $I_i \cap I_j$, for some $i \ne j$.

\begin{proposition}[conflict accounting]\label{prop:conflict}
Suppose $(\Lambda_1,\dots,\Lambda_\ell)$ is a free split selector tuple
satisfying $\Prob{\mathsf{Conf}} \le \eta$ together with the conclusion of
Conjecture~\ref{conj:main} for $Y[\Lambda]$ conditioned on
$\overline{\mathsf{Conf}}$. Then Theorem~\ref{thm:aibf} holds with $2\eta$
added to the bound of part~(1) and $\eta$ to that of part~(2), under any
merge rule.
\end{proposition}

\begin{proof}
Build $B$ as in Proposition~\ref{prop:sbf} but with
$L_i := \{(u,a_i(u)) : u \in I_i\}$. On $\overline{\mathsf{Conf}}$ the
merged list is $\bigcup_i L_i$ regardless of the rule, and the oracle at
$\main$ is $Y[\Lambda]$ conditioned on that event, so the argument runs
unchanged. Two probabilities that agree on an event of probability at
least $1-\eta$ differ by at most $\eta$, which for an
indistinguishability application enters the advantage doubled.
\end{proof}

\begin{remark}[the three rules, and which to aim for]\label{rem:rules}
Three regimes are worth distinguishing.
\begin{itemize}
\item \emph{Consistency.} $a_i = F|_{I_i}$. Then $\eta = 0$ by
  Proposition~\ref{prop:noconflict}, the merge rule is irrelevant, and the
  statement is symmetric in the stages. This is what
  Conjecture~\ref{conj:main} asks for, and what the single-source
  decomposition of~\cite{CDGS} already delivers.
\item \emph{Rare conflict.} A free tuple with $\Prob{\mathsf{Conf}} \le
  \eta$, costing $2\eta$. This is the right form for a proof that
  establishes a symmetric statement first and then bounds the overlaps,
  and it is strictly weaker than consistency.
\item \emph{Priority.} A free tuple with no bound on $\mathsf{Conf}$, the
  lemma being stated directly for the priority-merged source. No extra
  term is incurred, but the price is real and easy to overlook: on
  $\mathsf{Conf}$ the merged oracle contradicts some $L_j$, hence
  contradicts the advice $z_j$ was computed against, and it is precisely
  such contradictions between advice and oracle that an online adversary
  is built to detect. A lemma proved in this form must already account for
  that detection inside its own bound; nothing in
  Section~\ref{sec:theorem} does it for free. The statement also ceases to
  be symmetric under permutation of the stages.
\end{itemize}
Consistency is the one to aim for. It is no harder to ask for in any
construction where the selectors read the whole oracle
(Remark~\ref{rem:allofF}), and it is the only one of the three under which
the conclusion does not depend on an arbitrary tie-break. Some
applications are indifferent: Section~\ref{sec:ext} goes through under the
priority rule with no extra term, because its argument never compares the
oracle against the advice (Remark~\ref{rem:noconsistency}).
\end{remark}

\section{The role of unpredictable advice}\label{sec:unpred}

It is natural to require that no $z_i$ be guessable from the oracle alone,
since that is what forces the private coins to matter. The requirement is
recorded here, together with what it does and does not do.

\begin{definition}[unpredictable advice]\label{def:unpred}
The tuple $(A^{0}_1,\dots,A^{0}_\ell)$ is \emph{$\delta$-unpredictable} if
for every predictor $\mathsf{P}$ with unrestricted access to the whole
table, $\Prob{\mathsf{P}(F) = z_i} \le \delta$ for every $i$, the
probability being over $F$, the coins of $A^{0}_i$ and those of
$\mathsf{P}$. Equivalently
$\Exp_{F}\bigl[\max_z \Prob{z_i = z \mid F}\bigr] \le \delta$.
\end{definition}

\begin{remark}[the hypothesis is not needed, and marks the hard case]
\label{rem:det}
Conjecture~\ref{conj:main} and Theorem~\ref{thm:aibf} are stated without
it, deliberately. At the opposite extreme, if each $z_i$ is a
deterministic function of the table --- the case $\delta = 1$ --- the
theorem is immediate from~\cite{CDGS} and needs no conjecture: stage
$B^{0}_1$, which reads $F$, recomputes $z_j = f_j(F)$ for every $j \ge 2$
by itself, applies the single-source selector of~\cite[Lem.~1]{CDGS} to
$z_1\|\cdots\|z_\ell$ with budget $P_1 := P$, and presets it; every other
stage outputs $z_j = f_j(F)$ and presets nothing. Small $\delta$ therefore
\emph{strengthens} the challenge rather than easing it: it is exactly the
regime in which no stage can reconstruct another's advice. Adding
Definition~\ref{def:unpred} as a hypothesis would weaken
Theorem~\ref{thm:aibf} without shortening any step of its proof.
Section~\ref{sec:ext} supplies its own hypothesis for its own reasons, and
it is about the \emph{points} the advice encodes
(Definition~\ref{def:nu}), not about the advice itself.
\end{remark}

The bookkeeping the single-source proof rests on is unaffected by
randomized leakage, which is what confines the difficulty to splitness.

\begin{lemma}[entropy bookkeeping for randomized leakage]\label{lem:randleak}
Let $X \sample [M]^{N}$ and let $Z$ be any random variable on $\bits^{S}$
jointly distributed with $X$; in particular $Z$ may be $f(X;r)$ for a
randomized $f$. Put $S_z := N\log M - \Hmin(X \mid Z = z)$. Then
\[
\begin{aligned}
  \Exp_{z \sample Z}\bigl[S_z\bigr] &\;\le\; H(Z) \;\le\; S , \\
  \Prob{S_Z > S + \log \gamma^{-1}} &\;\le\; \gamma
  \quad\text{for every } \gamma \in (0,1).
\end{aligned}
\]
\end{lemma}

\begin{proof}
Since $X$ is uniform,
\[
  \Prob{X = x \mid Z = z} \;=\;
  \frac{M^{-N}\,\Prob{Z = z \mid X = x}}{\Prob{Z = z}}
\]
for every $z$ in the support of $Z$, so taking the maximum over $x$,
\begin{equation}\label{eq:Sz}
  S_z \;=\; \log \max_{x} \Prob{Z = z \mid X = x} \;-\; \log \Prob{Z = z}.
\end{equation}
Each first term is at most $0$, so averaging \eqref{eq:Sz} over
$z \sample Z$ gives $\Exp_z[S_z] \le -\sum_z \Prob{Z = z}\log\Prob{Z = z}
= H(Z) \le S$, the last step because $Z$ takes at most $2^{S}$ values. For
the tail, \eqref{eq:Sz} rearranges to $\Prob{Z = z} = 2^{-S_z}\max_x
\Prob{Z = z \mid X = x} \le 2^{-S_z}$, whence
\begin{align*}
  \Prob{S_Z > S + \log\gamma^{-1}}
  &\;=\; \sum_{z\,:\,S_z > S + \log\gamma^{-1}} \Prob{Z = z} \\
  &\;\le\; 2^{S}\cdot 2^{-(S + \log\gamma^{-1})} \;=\; \gamma . \qedhere
\end{align*}
\end{proof}

\begin{remark}[deterministic leakage as the equality case]\label{rem:claim4}
For deterministic $f$ one has $\max_x \Prob{Z = z \mid X = x} = 1$, so
\eqref{eq:Sz} reduces to $\Prob{Z = z} = 2^{-S_z}$ and both parts of
Lemma~\ref{lem:randleak} become Claim~4 of~\cite{CDGS}. The lemma is the
statement that randomization can only help, and it is why
Remark~\ref{rem:reduces} may assert that the single-source case of
Conjecture~\ref{conj:main} is known.
\end{remark}

\begin{remark}[a refinement that fails]\label{rem:effective}
It is tempting to expect $\delta$-unpredictability to buy a smaller
budget, on the ground that private coins dilute the advice: if $z$ costs
$S$ bits of which $\log(1/\delta)$ are the predictor's uncertainty, the
\emph{effective} leakage should be $S - \log(1/\delta)$, and $P_i$ should
scale with that rather than with $S_i$. The bookkeeping does not support
this. Let $X \sample [4]$ and let $Z$ be one bit with
$\Prob{Z = 0 \mid X = 1} = 1$, $\Prob{Z = 1 \mid X = 2} = 1$ and
$\Prob{Z = 0 \mid X = 3} = \Prob{Z = 0 \mid X = 4} = \tfrac12$. Then
$\Prob{Z = 0} = \Prob{Z = 1} = \tfrac12$ and
$\max_x \Prob{X = x \mid Z = z} = \tfrac12$ for both $z$, so $S_z = 1$ for
both and $\Exp_z[S_z] = 1 = S$, while
$\delta = \tfrac14(1+1+\tfrac12+\tfrac12) = \tfrac34$ and
$S - \log(1/\delta) = 1 - \log\tfrac43 \approx 0.585 < 1$. The average
min-entropy deficiency is not bounded by the effective leakage, so any
refinement along these lines needs a different functional of $\delta$, or
a different argument.
\end{remark}

\section{Multi-source extraction}\label{sec:ext}

Now let the domain factor as $[N] = [N_1]\times\cdots\times[N_\ell]$, so
$N = \prod_i N_i$. Stage $i$ outputs
$z_i = (x_i,\ell_i) \in [N_i]\times\bits^{\sigma_i}$, so that
$S_i = \lceil\log N_i\rceil + \sigma_i$ and $\sigma := \sum_i \sigma_i$;
write $\vx := (x_1,\dots,x_\ell)$ and $\vl := (\ell_1,\dots,\ell_\ell)$,
and take $\pi(\vz) := \vl$. Such a tuple of stages is a \emph{split source
tuple}: each reads the whole oracle, they run on mutually independent
coins and never communicate. The extractor is the oracle itself at the
tuple of points, $F(\vx)$, and the observer sees $\vl$ and a challenge
value but never $\vx$.

\subsection{Guessing probabilities}\label{sec:delta}

\begin{definition}[single-coordinate guessing probabilities]
\label{def:delta}
For $i \in [\ell]$ put
\[
  \delta^{\ast}_i \;:=\; \sup_{\mathsf{P}} \Prob{\mathsf{P}(F) = x_i},
  \qquad
  \delta_i \;:=\; \sup_{\mathsf{P}} \Prob{\mathsf{P}(F,\ell_i) = x_i},
\]
the suprema over all predictors, the probabilities over $F \sample \Fun$,
the coins $\vr$ and those of $\mathsf{P}$. Both predictors hold the
\emph{entire} table, so neither is granted or denied oracle queries; only
the leakage distinguishes them. Write
$\delta^{\ast} := \max_i \delta^{\ast}_i$ and $\delta := \max_i \delta_i$.
\end{definition}

The parameter $\delta^{\ast}_i$ is the hypothesis in its informal reading:
no predictor given access to the full random oracle can guess a source's
output.

\begin{lemma}[the leakage costs at most its length]\label{lem:leak}
$\delta^{\ast}_i \le \delta_i \le 2^{\sigma_i}\delta^{\ast}_i$, and
$\delta_i = \delta^{\ast}_i$ when $\sigma_i = 0$.
\end{lemma}

\begin{proof}
The first inequality holds because a predictor may ignore $\ell_i$. For
the second, write
\[
  \delta_i \;=\; \Exp_{F}\Bigl[\, \sum_{\ell} \Prob{\ell_i = \ell \mid F}
  \;\max_u \Prob{x_i = u \mid F, \ell_i = \ell} \Bigr] .
\]
Each summand is at most $\max_u \Prob{x_i = u \mid F}$, since
\begin{align*}
  \Prob{\ell_i = \ell \mid F}\,\Prob{x_i = u \mid F,\ell_i=\ell}
  &\;=\; \Prob{x_i = u,\ \ell_i = \ell \mid F} \\
  &\;\le\; \Prob{x_i = u \mid F} ,
\end{align*}
and there are at most $2^{\sigma_i}$ summands, while $\delta^{\ast}_i$ is
the expectation over $F$ of that same maximum.
\end{proof}

\begin{lemma}[the siblings' views are free]\label{lem:sibling}
For every $i$ and every predictor $\mathsf{P}$, writing
$W_{-i} := (z_j,r_j)_{j \ne i}$,
\[
  \Prob{\mathsf{P}(F,\ell_i,W_{-i}) = x_i} \le \delta_i ,
  \qquad
  \Prob{\mathsf{P}(F,W_{-i}) = x_i} \le \delta^{\ast}_i .
\]
\end{lemma}

\begin{proof}
Conditioned on $F$, each pair $(z_j,r_j)$ is a function of $r_j$ alone and
the $r_j$ are mutually independent, so $(z_i,r_i) \perp W_{-i} \mid F$.
Conditioning further on $\ell_i$, a function of $(F,r_i)$, preserves this,
so $x_i \perp W_{-i} \mid (F,\ell_i)$. A predictor's optimal success
probability against a target is unchanged by extra inputs conditionally
independent of the target given what it already holds, since the optimum
is $\Exp[\max_u \Prob{x_i = u \mid \text{inputs}}]$ and the conditional
law does not move.
\end{proof}

What the proof actually charges for is not a single coordinate but a
tuple: a selector must complete the challenge point, and an observer must
name all of it.

\begin{definition}[tuple guessing probabilities]\label{def:nu}
\begin{align*}
  \nu &\;:=\; \max_{i}\ \sup_{\mathsf{P}}
  \Prob{\mathsf{P}(F,z_i,r_i) = (x_j)_{j \ne i}} , \\
  \nu^{+} &\;:=\; \sup_{\mathsf{P}} \Prob{\mathsf{P}(F,\vl) = \vx} .
\end{align*}
The first is the chance of completing the challenge point from one
source's \emph{entire} view; the second, of naming all of it from the
oracle and every leakage.
\end{definition}

\begin{lemma}[what the tuple parameters are worth]\label{lem:nu}
$\nu \le \delta^{\ast}_{(2)}$, the second smallest of the
$\delta^{\ast}_j$, and $\nu^{+} \le \min_j \delta_j$. At $\ell = 2$ these
read $\nu \le \delta^{\ast}$ and $\nu^{+} \le \min_j \delta_j$. The
optimal predictor in the first supremum achieves
$\Exp_F\bigl[\prod_{j \ne i}\max_u \Prob{x_j = u \mid F}\bigr]$.
\end{lemma}

\begin{proof}
A predictor may discard all but the $j$-th coordinate of its output, so
the $i$-th supremum is at most $\min_{j \ne i}\delta^{\ast}_j$ by the
second part of Lemma~\ref{lem:sibling}; maximising over $i$ gives the
second smallest, since $\min_{j \ne i}\delta^{\ast}_j$ is the smallest for
every $i$ except the argmin and the second smallest for that one.
Likewise $\nu^{+} \le \delta_j$ for each $j$ by the first part, whose
input vector contains $(F,\vl)$. For the last claim, conditioned on $F$
the variables $x_j$, $j \ne i$, are mutually independent and independent
of $(z_i,r_i)$, so the optimal predictor names each coordinate's
conditional mode.
\end{proof}

\begin{remark}[$\nu$ is not the product of the $\delta^{\ast}_j$]
\label{rem:notproduct}
The expression in Lemma~\ref{lem:nu} is an expectation of a product, where
$\prod_{j \ne i}\delta^{\ast}_j$ is a product of expectations, and the two
are not ordered. If the sources are jointly predictable on half the tables
and jointly unpredictable on the other half, the expectation of the
product is about $\tfrac12$ while the product of expectations is
$2^{-(\ell-1)}$. So $\nu = \Theta(\delta^{\ell-1})$ holds when the
per-table guessing profiles are uncorrelated across the oracle, and $\nu$
can be as large as $\delta^{\ast}_{(2)}$ when they are not.
Remark~\ref{rem:ell} takes this up.
\end{remark}

\subsection{The game, and why the model needs both restrictions}
\label{sec:game}

\begin{definition}[the extraction application and its experiments]
\label{def:games}
$G^{\Ext}$ is the $\pi$-restricted indistinguishability application, for
$\pi(\vz) = \vl$, in which $\mathsf{C}(\vz)$ draws $b \sample \bits$, sets
$y := \mathcal{O}.\main(\vx)$ if $b = 1$ and $y \sample [M]$ if $b = 0$,
sends $y$ to the online stage, receives a bit $b'$ and outputs
$\llbracket b' = b\rrbracket$. The online stage is an \emph{observer} $D$
making at most $q$ queries, so $\Tcomb = q+1 =: q^{+}$.

On one probability space --- $X \sample \Fun$ and $\vr$ drawn once,
$z_i := f_i(X;r_i)$, never resampled --- and for a split selector tuple
with $Y[\Lambda]$, $I$ as in Definition~\ref{def:selector}, define
\begin{itemize}
\item $\Real$: $y := X(\vx)$; output $D^{X}(\vl,y)$;
\item $\Real_0$: as $\Real$ but $y \sample [M]$;
\item $\Dec$: $y := Y[\Lambda](\vx)$; output $D^{Y[\Lambda]}(\vl,y)$;
\item $\Dec_0$: as $\Dec$ but $y \sample [M]$.
\end{itemize}
The \emph{extraction advantage} is
$\kappa(q) := \sup_{D} |\Prob{\Real = 1} - \Prob{\Real_0 = 1}|$, the
supremum over $q$-query observers; it refers to $\Real$ and $\Real_0$
only, so it does not depend on $\Lambda$. Writing
$A := (\Sigma_1,\dots,\Sigma_\ell,D)$ for the split attacker whose
preprocessing is the sources and whose online stage is $D$, and $B$ for
the split bit-fixing attacker of Proposition~\ref{prop:sbf}, the identity
$2|\Succ - \tfrac12| = |\Prob{b'{=}1 \mid b{=}1} - \Prob{b'{=}1 \mid
b{=}0}|$ gives
\begin{equation}\label{eq:kappa-is-adv}
\begin{aligned}
  \Adv_{G^{\Ext},\SAI}(A)
  &= \bigl|\Prob{\Real = 1} - \Prob{\Real_0 = 1}\bigr| , \\
  \Adv_{G^{\Ext},\SBF}(B)
  &= \bigl|\Prob{\Dec = 1} - \Prob{\Dec_0 = 1}\bigr| .
\end{aligned}
\end{equation}
\end{definition}

\begin{remark}[why the online stage must be restricted]\label{rem:whypi}
The challenger must know $\vx$ to evaluate the extractor and the observer
must not, which is why Definition~\ref{def:app} carries $\pi$. Nothing
would be gained by taking $\pi$ the identity and hoping the restriction is
harmless: an online stage given $\vx$ queries the oracle there and
compares, winning with probability $1-1/M$ in the $\SAI$-model \emph{and}
in the $\SBF$-model, where the challenge is again
$\mathcal{O}.\main(\vx)$, so the transfer would be true and empty. The
content of extraction is precisely the gap between what the challenger
must know and what the adversary may know.
\end{remark}

\begin{remark}[the split bit-fixing model may not be collapsed]
\label{rem:nomerge}
$G^{\Ext}$ is \emph{insecure} in the ordinary bit-fixing model at every
parameter setting: a single preprocessing stage that produces the advice
knows every coordinate of $\vx$, so it may fix $\vx$ and preset the oracle
there to a constant $v$ written into its own description; the online
stage, equally non-uniform, holds $v$ and outputs
$\llbracket y = v\rrbracket$, distinguishing with probability $1-1/M$ and
no queries at all. No leakage is needed for this, so the failure covers
Corollary~\ref{cor:noleak} as well, and no restriction on the advice
distribution repairs it, since the offending stage's own view determines
$\vx$. So Proposition~\ref{prop:merge} and hence
Theorem~\ref{thm:aibf} are unavailable here: the derivation stops at
Proposition~\ref{prop:sbf} and proves its bit-fixing bound directly. That
bound is short, and it is where the independence of the stages is spent.
Note also, per Remark~\ref{rem:derandcost}, that
Proposition~\ref{prop:derand} must not be applied: hardwiring a table
chosen to maximise success leaves $\nu$ unconstrained, and
Lemmas~\ref{lem:hit}--\ref{lem:bfext} would fail of the derandomized
attacker. Nothing here needs $B$ to be admissible.
\end{remark}

\subsection{The three bounds in the bit-fixing world}\label{sec:bf}

Throughout, $(\Lambda_1,\dots,\Lambda_\ell)$ is the selector tuple
Conjecture~\ref{conj:main} supplies for
$f_1,\dots,f_\ell,P_1,\dots,P_\ell,\gamma$, and $I_i := I_i(X,z_i,r_i)$,
$I = \bigcup_i I_i$ are its fixed sets on the probability space of
Definition~\ref{def:games}.

\begin{lemma}[the fixed set misses the challenge]\label{lem:hit}
$\Prob{\vx \in I} \le \sum_i P_i\nu = P\nu$.
\end{lemma}

\begin{proof}
Fix $i$. Let $\mathsf{P}$ receive $(X,z_i,r_i)$, compute
$(I_i,a_i) := \Lambda_i(X,z_i,r_i)$ --- it can, since $\Lambda_i$ depends
only on $f_1,\dots,f_\ell,P_1,\dots,P_\ell,\gamma$ and those three inputs
--- and output the coordinates $j \ne i$ of a uniformly chosen element of
$I_i$, or an arbitrary tuple if $I_i$ is empty. If $\vx \in I_i$ then that
element is $\vx$ with probability at least $1/|I_i| \ge 1/P_i$, and its
coordinates $j \ne i$ are then $(x_j)_{j \ne i}$, so
$\Prob{\mathsf{P}(X,z_i,r_i) = (x_j)_{j\ne i}} \ge \Prob{\vx \in I_i}/P_i$,
which is at most $\nu$ by Definition~\ref{def:nu}. Sum over $i$.
\end{proof}

\begin{remark}[this is where splitness is spent]\label{rem:onesource}
Lemma~\ref{lem:hit} is the only step that uses the stages' independence in
an essential way, and it uses it in a form that has no single-source
analogue. The selector $\Lambda_i$ is allowed to read $z_i$, hence $x_i$,
so it knows one coordinate of the challenge point exactly; what it cannot
do is complete the tuple, and $P_i\nu$ is the cost of the completion. At
$\ell = 1$ the complementary tuple is empty, so $\nu = 1$ and the lemma is
vacuous --- the right answer, since extraction from one source is false: a
source that queries $F$ at its own output point and leaks one bit of the
answer defeats every admissible family.
\end{remark}

\begin{lemma}[the observer misses the challenge]\label{lem:query}
Let $Q$ be the set of points the observer queries. Then
$\Prob{\vx \in Q} \le q\nu^{+}$ in $\Real_0$, and
$\Prob{\vx \in Q} \le q\nu^{+} + \Delta(q) + \gamma$ in $\Dec_0$. Both
hold trivially at $q = 0$, where $Q$ is empty.
\end{lemma}

\begin{proof}
In $\Real_0$: let $\mathsf{P}$ receive $(X,\vl)$, draw $y \sample [M]$ and
the observer's coins, run $D^{X}(\vl,y)$ recording its queries, and output
a uniformly chosen one. It can run $D$ faithfully because the oracle in
$\Real_0$ \emph{is} $X$, which it holds in full, and because the challenge
there is uniform and independent of everything else, so its own draw of
$y$ has the right law; its execution therefore has the law of the one in
$\Real_0$ conditioned on $(X,\vx,\vz)$. As $|Q| \le q$,
$\Prob{\mathsf{P}(X,\vl) = \vx} \ge \Prob{\vx \in Q}/q$, which is at most
$\nu^{+}$.

In $\Dec_0$ the oracle is $Y[\Lambda]$, whose fixed set depends on $\vr$,
and a predictor holding only $(X,\vl)$ cannot reconstruct it; the
conjecture bridges the gap. Let $D'$ take the tuple $\vz$, parse it as
$(\vx,\vl)$, draw $y \sample [M]$ and the observer's coins, run
$D(\vl,y)$ forwarding its at most $q$ oracle queries, and output $1$ if
one of those queries equals $\vx$. It is a $q$-query distinguisher taking
a tuple of advice strings, which is what Conjecture~\ref{conj:main}
quantifies over --- that it reads the private part of the advice is
permitted there, and is the same latitude Proposition~\ref{prop:sbf} uses.
By construction $\Prob{D'^{X}(\vz)=1}$ and
$\Prob{D'^{Y[\Lambda]}(\vz)=1}$ are the two probabilities in question, and
the first bound of the conjecture at $T := q$, with the \emph{same}
selector tuple --- legitimate because it is fixed before $T$ and before
the distinguisher (Remark~\ref{rem:order}) --- separates them by at most
$\Delta(q)+\gamma$.
\end{proof}

\begin{lemma}[extraction in the split bit-fixing model]\label{lem:bfext}
Writing $Q$ for the set of points the observer queries in $\Dec_0$,
\[
  \bigl|\Prob{\Dec = 1} - \Prob{\Dec_0 = 1}\bigr|
  \;\le\; \Prob{\vx \in I} + \Prob{\vx \in Q} .
\]
\end{lemma}

\begin{proof}
Condition on $(X,\vz,\vr)$, which fixes $\vx$, $\vl$, the fixed set $I$
and the fixed values; the two experiments differ only in the challenge. On
the event $\vx \in I$ bound the difference by $1$; this contributes
$\Prob{\vx \in I}$. Assume $\vx \notin I$. By
Definition~\ref{def:selector} the values of $Y[\Lambda]$ off $I$ are
uniform on $[M]$ and independent across points and of the conditioning, so
$Y[\Lambda](\vx)$ is uniform and independent of both
$(Y[\Lambda](u))_{u \ne \vx}$ and $\vl$.

Couple as follows. Draw $Y[\Lambda]$ and an independent $U \sample [M]$,
and let $\widetilde{Y}$ agree with $Y[\Lambda]$ except that
$\widetilde{Y}(\vx) := U$. Since the value at $\vx$ is uniform and
independent of the rest, $\widetilde{Y}$ has the same law as $Y[\Lambda]$
given the conditioning, so running the observer on $(\widetilde{Y},U)$
reproduces $\Dec$, whose challenge is the oracle's own value at $\vx$, and
running it on $(Y[\Lambda],U)$ reproduces $\Dec_0$, whose challenge is an
independent uniform value. The two executions receive the same input
$(\vl,U)$ and the same coins, and the two oracles differ only at $\vx$;
hence they proceed identically up to and including the step before a first
query at $\vx$, so that event has the same probability in both, and on its
complement the two outputs agree. Bounding it in $\Dec_0$ and averaging
over the conditioning gives the claim.
\end{proof}

\begin{remark}[consistency is not needed]\label{rem:noconsistency}
Lemma~\ref{lem:bfext} uses only that $|I| \le P$ and that $Y[\Lambda]$ is
uniform and independent off $I$; the fixed \emph{values} never appear. The
same is true of Lemmas~\ref{lem:hit} and~\ref{lem:query}, which speak of
the fixed sets alone. Everything in this section therefore follows from
the weakening of Conjecture~\ref{conj:main} in which the selector tuple is
\emph{free} (Section~\ref{sec:conflict}), with no bound whatever on the
conflict probability and under any merge rule --- the priority regime of
Remark~\ref{rem:rules}, whose hidden cost is the possibility that the
merged oracle contradicts the advice. That cost is real for a general
application; for $G^{\Ext}$ it is not incurred, because the observer's
only handle on the advice is $\vl$ and the argument never compares the
oracle against it.
\end{remark}

\subsection{The extraction theorem}\label{sec:extthm}

\begin{theorem}[multi-source extraction from parallel presampling]
\label{thm:ext}
Assume Conjecture~\ref{conj:main} with constant $c$. Let
$(\Sigma_1,\dots,\Sigma_\ell)$ be a split source tuple with advice lengths
$S_1,\dots,S_\ell$ and tuple guessing probabilities $\nu, \nu^{+}$. Then
for every $q \in \mathbb{N}\cup\{0\}$, every $P_1,\dots,P_\ell$ and every
$\gamma \in (0,1)$,
\[
  \kappa(q) \;\le\; 3c\,q^{+}\sum_{i=1}^{\ell}\frac{S_i+\log\gamma^{-1}}{P_i}
  \;+\; P\nu \;+\; q\nu^{+} \;+\; 3\gamma ,
\]
and, optimising, with $C_\ell := 2\sqrt{3c\,\ell(\ell+1)} + \ell + 3$,
\[
  \kappa(q) \;\le\; C_\ell\sqrt{S\,q^{+}\,\nu} \;+\; q\nu^{+}
  \;=\; O\!\left(\ell\sqrt{S\,q^{+}\,\nu}\right)
\]
whenever the bound is below one.
\end{theorem}

\begin{proof}
Fix a $q$-query observer $D$, so that $\Tcomb = q^{+}$, and let $A$ and
$B$ be as in Definition~\ref{def:games}. Applying the first bound of
Proposition~\ref{prop:sbf} to the indistinguishability application
$G^{\Ext}$, as in the proof of Theorem~\ref{thm:aibf} but without its
appeal to Proposition~\ref{prop:merge}, and using
\eqref{eq:kappa-is-adv},
\[
\begin{aligned}
  \bigl|\Prob{\Real{=}1}-\Prob{\Real_0{=}1}\bigr|
  \;\le\;& \bigl|\Prob{\Dec{=}1}-\Prob{\Dec_0{=}1}\bigr| \\
  &+\; 2\Delta(q^{+}) + 2\gamma ,
\end{aligned}
\]
and Lemmas~\ref{lem:bfext},~\ref{lem:hit} and~\ref{lem:query} bound the
first term on the right by $P\nu + q\nu^{+} + \Delta(q) + \gamma$. Since
$\Delta$ is increasing in $T$ and $q \le q^{+}$, adding gives
$\kappa(q) \le 3\Delta(q^{+}) + P\nu + q\nu^{+} + 3\gamma$, which is the
first display.

For the second, choose $P_i := \lceil\sqrt{3cq^{+}A_i/\nu}\,\rceil$, which
is legitimate even though it depends on $q$: the conjecture supplies a
tuple for \emph{every} $(P_1,\dots,P_\ell,\gamma)$ and we instantiate it
at one of them. What Remark~\ref{rem:order} forbids is a tuple depending
on $T$ or on the distinguisher, and the two applications made above --- at
$T = q^{+}$ in Proposition~\ref{prop:sbf} and at $T = q$ in
Lemma~\ref{lem:query} --- use the same tuple. From
$P_i \ge \sqrt{3cq^{+}A_i/\nu}$,
\[
  \frac{3cq^{+}A_i}{P_i} \le \sqrt{3cq^{+}A_i\nu},
  \qquad
  P_i\nu \le \sqrt{3cq^{+}A_i\nu} + \nu ,
\]
so the first two terms of the first display total at most
\[
  2\sqrt{3cq^{+}\nu}\,\sum_i \sqrt{A_i} \;+\; \ell\nu ,
\]
and $\sum_i\sqrt{A_i} \le \sqrt{\ell A}$ by Cauchy--Schwarz, where
$A := \sum_i A_i = S + \ell\log\gamma^{-1}$.

Now put $\gamma := N^{-1} = (\prod_i N_i)^{-1}$. Then
$\log\gamma^{-1} = \sum_i \log N_i \le S$, so $A \le (\ell+1)S$; and
$\gamma \le \nu$, because a predictor guessing the complementary tuple
uniformly already achieves $\prod_{j\ne i}N_j^{-1} \ge \gamma$.
Collecting,
\[
  \kappa(q) \;\le\; 2\sqrt{3c\,\ell(\ell+1)\,Sq^{+}\nu}
  + (\ell+3)\nu + q\nu^{+} ,
\]
and $\nu \le \sqrt{Sq^{+}\nu}$ because $\nu \le 1 \le Sq^{+}$.
\end{proof}

\begin{corollary}[two sources]\label{cor:two}
At $\ell = 2$, Lemma~\ref{lem:nu} gives $\nu \le \delta^{\ast}$ and
$\nu^{+} \le \delta$, so
$\kappa(q) \le \bigl(6\sqrt{2c}+5\bigr)\sqrt{S q^{+}\delta^{\ast}} +
q\delta$.
\end{corollary}

\begin{corollary}[the leakage-free case]\label{cor:noleak}
If the sources output no public leakage, so that every $\sigma_i = 0$ and
the observer sees only the challenge value, then
$\delta_i = \delta^{\ast}_i$ by Lemma~\ref{lem:leak},
$\nu^{+} \le \min_j\delta^{\ast}_j$ and $\nu \le \delta^{\ast}_{(2)}$,
whence
\[
  \kappa(q) \;\le\; C_\ell
  \sqrt{\textstyle\bigl(\sum_i \lceil\log N_i\rceil\bigr)\,q^{+}\,
  \delta^{\ast}_{(2)}} \;+\; q\min_j \delta^{\ast}_j .
\]
This is the statement in its plainest form: if no $x_i$ can be guessed
from the entire random oracle, then $F(x_1,\dots,x_\ell)$ is
indistinguishable from uniform against $q$ queries.
\end{corollary}

\begin{remark}[where each hypothesis is used]\label{rem:uses}
The main term carries $\nu$, controlled by unpredictability given the
oracle \emph{alone}; the leakage-conditioned parameter enters only through
$\nu^{+}$ in the additive $q\nu^{+}$, which is unavoidable, since an
observer holding $\vl$ and $q$ queries can spend them on its best guesses
at $\vx$. By Lemma~\ref{lem:leak} the two can differ by $2^{\sigma}$, so
the distinction is not cosmetic once the leakage is long. Of the
conjecture, only the first (additive) bound is used; the multiplicative
one has no role, extraction being an indistinguishability application. Of
the selector tuple, only the budgets $|I_i| \le P_i$, the splitness of the
selectors (Lemma~\ref{lem:hit}), and the quantifier order, which lets one
tuple serve two distinguishers at two different query budgets.
\end{remark}

\section{Discussion}\label{sec:discussion}

\begin{remark}[the statement is known; the route is what is new]
\label{rem:priorart}
Theorem~\ref{thm:ext} is not a new statement. \cite[Thm.~3]{CFHS} proves
multi-source extraction for the monolithic random oracle
\emph{unconditionally}, for $\ell \ge 2$ arbitrary sources with unbounded
oracle access that are $\delta$-unpredictable in exactly the sense of
Definition~\ref{def:delta} --- their predictor reads the whole table and
holds the source's own leakage --- against a $q_D$-query distinguisher, by
a compression argument in the manner of~\cite{DGK} combined with the flat
decomposition of high-min-entropy sources. Their bound is
\[
  O\!\left(\sqrt{\ell^{\,\ell+1}\log M\,(\sigma+\ell N_1)\,q_D\,
  \delta^{\ell}}\right) ,
\]
where $N_1$ is the size of a single coordinate's domain, against
$O(\ell\sqrt{Sq^{+}\nu})$ here, and the two are incomparable. Ours is
logarithmic in the domain size where theirs carries $\sqrt{N_1}$, and
carries no $\ell^{\ell}$; theirs is linear in $\delta$ at $\ell = 2$ and
improves as $\delta^{\ell/2}$, where ours stays a square root of a single
guessing probability at every $\ell$ (Remark~\ref{rem:ell}).

At $\ell = 2$ the comparison favours the bound here almost everywhere.
Ours is the smaller once $Sq^{+}\delta^{\ast} \lesssim N_1\log M\,q\delta^{2}$,
that is, once $\delta \gtrsim S/(N_1\log M)$, and since $\delta \ge 1/N_1$
always, that threshold sits at or below the floor whenever
$\log M \gtrsim S$: for short leakage and an output range no smaller than
the advice, there is no $\delta$ at which theirs is better. Theirs wins
where the leakage is long or $M$ is small, and, for $\ell \ge 3$, over a
range that widens with $\ell$.

What is derived here is therefore the \emph{route}, not the conclusion.
\cite[\S1.4]{CFHS} records that decomposition ``turns out not to be the
case, at least when these techniques are applied directly in parallel'',
gives the counterexample of Remark~\ref{rem:separate}, and conjectures
that ``parallel decomposition is possible under appropriate
restrictions''. Theorem~\ref{thm:ext} is the first half of what that would
buy: it confirms that Conjecture~\ref{conj:main}, if true, does deliver
the extraction bound, and it exhibits the places where splitness has to be
spent to get it (Lemma~\ref{lem:hit} and Remark~\ref{rem:nomerge}). The
reason~\cite{CFHS} wants the route in the first place is the \emph{second}
half --- non-monolithic constructions, where compression stalls --- and
that is exactly what Remark~\ref{rem:construction} does not cover. Until
it is covered, the conjecture buys, for extraction, no unconditional
statement that is not already in~\cite{CFHS}.
\end{remark}

\begin{remark}[the route does not profit from extra sources]\label{rem:ell}
The main term is $\sqrt{Sq^{+}\nu}$ at every $\ell$: a square root, of a
single guessing probability, no matter how many sources there are, with
only the $O(\ell)$ in $C_\ell$ recording their number.
\cite[Thm.~3]{CFHS} by contrast reaches $\delta^{\ell/2}$, improving with
each additional source. The gap is not an artefact of the bookkeeping but
of what the two proofs charge for. Presampling charges twice, once for the
fixed set and once for the observer's queries, and each charge is a single
guessing event; adding sources makes that event harder only through $\nu$,
and by Remark~\ref{rem:notproduct} that is $\Theta(\delta^{\ell-1})$ only
when the per-table guessing profiles are uncorrelated across the oracle.
Even in that best case the bound reads
$O(\ell\,\delta^{(\ell-1)/2}\sqrt{Sq^{+}})$, still one half-power of
$\delta$ behind~\cite[Thm.~3]{CFHS}. Whether that half-power is the price
of the route or of this particular argument, we do not know.
\end{remark}

\begin{remark}[non-monolithic extractors]\label{rem:construction}
The motivation in~\cite[\S1.4]{CFHS} for wanting a decomposition route at
all is that compression stalls for constructions such as
Merkle--Damg\aa{}rd and sponge, where the extractor is not a single oracle
call. Theorem~\ref{thm:ext} is proved for the monolithic extractor
$F(\vx)$ and does not by itself cover those. The obstruction is localised:
for an extractor $\mathsf{C}^{F}(\vx)$ making $t$ queries one would take
$\Tcomb = q+t$, and Lemma~\ref{lem:bfext} would need the whole of
$\mathsf{C}$'s query set, not merely $\vx$, to avoid the fixed set and the
observer's queries. The first blocks queried are unpredictable for the
same reason $\vx$ is; the later ones are determined by the answers to the
earlier ones, and bounding their collision with $I$ needs an argument this
document does not make. We record this as the natural next step rather
than as a corollary.
\end{remark}

\begin{remark}[the companion mixture statement]\label{rem:mixture}
A different decomposition statement for the same setting is worth
recording, both because the Conjura page carried it through revision~3 and
because it is what the campaign in that repository attacks. There the
parties output \emph{points} $x_i$ together with leakage $z_i$; the object
asked for is a single \emph{joint} $P$-mixture $Y_{f,\zeta}$, indexed by
the oracle and all the leakages but \emph{not} by the points, every
component consistent with $f$; and the observer is
\emph{challenge-oblivious} by fiat, seeing the leakage and the value
$H(\vx)$ but never $\vx$. Its conjectured bound, for
$\sigma' := \sigma + \ell\log N_1$, is
$c(\sigma'+\log\gamma^{-1})q^{+}/P + C\sqrt{\sigma'q^{+}\delta} + \gamma$,
and its consequence is extraction of the same shape as
Theorem~\ref{thm:ext}.

Neither is known to imply the other, and they differ in three places at
once. The mixture there is joint, chosen with every leakage in hand, where
the selectors here are split; but it is forbidden the points, where
$\Lambda_i$ may read its own $x_i$ --- which is why Lemma~\ref{lem:hit}
must spend splitness to pay $P\nu$ where the mixture form pays $P\delta$
by that prohibition instead. And its observer is challenge-oblivious by
definition, where here the restriction belongs to the application
(Definition~\ref{def:app}) and the conjecture's own distinguisher reads
everything. What is proved of each also differs: nothing is proved of
Conjecture~\ref{conj:main}, whereas at $\ell = 2$ the mixture form's
query-free case is proved and tight and its one-query case is proved up to
polylogarithmic factors under an extra hypothesis. Those results attach to
the mixture form, not to Conjecture~\ref{conj:main}.
\end{remark}

\begin{remark}[what is not claimed]\label{rem:notclaimed}
Everything after Section~\ref{sec:lemma} is conditional on
Conjecture~\ref{conj:main}, which is open; the deterministic-preprocessing
case of that conjecture is immediate from~\cite{CDGS} but is also, by
Remark~\ref{rem:det}, the case in which the hypothesis of
Section~\ref{sec:ext} is vacuous, since a source whose output is a
deterministic function of the oracle has $\delta^{\ast}_i = 1$, and then
$\nu = 1$ as soon as two sources are deterministic. The interesting regime
for extraction is therefore exactly the regime in which the conjecture is
open, and Theorem~\ref{thm:ext} should be read as locating the extraction
question inside the presampling one, not as settling it.
\end{remark}

\section{Bibliography}

\begin{conjurabibliography}{9}

\bibitem[CDGS]{CDGS}
S. Coretti, Y. Dodis, S. Guo, J. Steinberger.
\emph{Random Oracles and Non-Uniformity}.
EUROCRYPT 2018; ePrint 2017/937.

\bibitem[CFHS]{CFHS}
S. Coretti, P. Farshim, P. Harasser, K. Southern.
\emph{Multi-Source Randomness Extraction and Generation in the
Random-Oracle Model}.
ITC 2025, LIPIcs 343, art. 10; ePrint 2025/1258.

\bibitem[DGK]{DGK}
Y. Dodis, S. Guo, J. Katz.
\emph{Fixing Cracks in the Concrete: Random Oracles with Auxiliary Input,
Revisited}.
EUROCRYPT 2017.

\bibitem[Unruh]{Unruh}
D. Unruh.
\emph{Random Oracles and Auxiliary Input}.
CRYPTO 2007.

\end{conjurabibliography}

\end{document}
